\section{Conclusion}
This work introduces the first orthonormal basis functions matched-filter framework for passive magnetic detection of DC current-carrying telecommunication cables. The approach offers a physically interpretable alternative to data-driven methods. \highlight{On both real lake data and controlled simulations, the proposed cable-specific OBF outperformed the multipolar OBF baselines (dipole and dipole--quadrupole) in localization accuracy, reducing the cable-crossing error to 0.6~m in the field experiment and 0.10~m in simulation, while remaining robust at low SNR.} \highlight{This result warrants an important distinction. At extremely low SNR, a matched detector built from the fewest basis functions is naturally favored, and the multipolar OBF (MOBF) baselines indeed achieve satisfactory detection in these tests. However good this detection performance may be, it is the localization capability, and in particular the isolation of local maxima, that the multipolar approximations fail to match at any SNR. This is a central concern of the present work, beyond detection alone.} Future work will investigate extensions to a gradiometer configuration and refined estimation of cable position and orientation. \highlight{A further direction concerns the straight-conductor assumption underlying the model: the real cable is only locally straight and neither perfectly stretched nor infinite. Although this effect was subtly noticed during the experiments, we did not develop a specific formulation or detailed analysis of it, and characterizing its impact on detection and localization remains an open problem.}